% =============================================================
% Timeline do SDDP — cada rank resolve um subconjunto distinto
% de cenários, alternando computação e E/S independente.
% Bibliotecas necessárias: arrows.meta, positioning
% =============================================================
\begin{figure}[H]
  \centering
  \definecolor{compcolor}{HTML}{66BB6A}
  \definecolor{iocolor}{HTML}{E57373}
  \begin{tikzpicture}[
    rankbg/.style={align=right, font=\small, anchor=east},
    compute/.style={draw, fill=compcolor, text=white, minimum height=0.6cm, anchor=west, font=\scriptsize\bfseries, inner sep=2pt},
    io/.style={draw, fill=iocolor, text=white, minimum height=0.6cm, anchor=west, font=\scriptsize\bfseries, inner sep=2pt},
    timeaxis/.style={-{Latex[length=2mm]}, thick}
  ]
    %---- rótulos de rank (à esquerda) ----
    \node[rankbg] at (-0.2, 0)    {rank 0};
    \node[rankbg] at (-0.2, -0.8) {rank 1};
    \node[rankbg] at (-0.2, -1.6) {rank 2};
    \node[rankbg] at (-0.2, -2.4) {rank 3};
    \node[font=\large, anchor=east] at (-0.2, -3.0) {$\vdots$};
    \node[rankbg] at (-0.2, -3.6) {rank $N$};

    %---- timelines: cada rank resolve um subconjunto distinto de cenários ----
    % rank 0: 1, 2, 3
    \node[compute, minimum width=3.5cm] at (0,    0) {Resolvendo cenário 1};
    \node[io,      minimum width=0.75cm] at (3.5,  0) {E/S};
    \node[compute, minimum width=3.5cm] at (4.25, 0) {Resolvendo cenário 2};
    \node[io,      minimum width=0.75cm] at (7.75, 0) {E/S};
    \node[compute, minimum width=3.5cm] at (8.5,  0) {Resolvendo cenário 3};
    \node[io,      minimum width=0.75cm] at (12.0, 0) {E/S};

    % rank 1: 4, 5, 6
    \node[compute, minimum width=3.8cm] at (0,    -0.8) {Resolvendo cenário 4};
    \node[io,      minimum width=0.75cm] at (3.8,  -0.8) {E/S};
    \node[compute, minimum width=3.4cm] at (4.55, -0.8) {Resolvendo cenário 5};
    \node[io,      minimum width=0.75cm] at (7.95, -0.8) {E/S};
    \node[compute, minimum width=3.5cm] at (8.7,  -0.8) {Resolvendo cenário 6};
    \node[io,      minimum width=0.75cm] at (12.2, -0.8) {E/S};

    % rank 2: 7, 8, 9
    \node[compute, minimum width=3.2cm] at (0,    -1.6) {Resolvendo cenário 7};
    \node[io,      minimum width=0.75cm] at (3.2,  -1.6) {E/S};
    \node[compute, minimum width=3.9cm] at (3.95, -1.6) {Resolvendo cenário 8};
    \node[io,      minimum width=0.75cm] at (7.85, -1.6) {E/S};
    \node[compute, minimum width=3.4cm] at (8.6,  -1.6) {Resolvendo cenário 9};
    \node[io,      minimum width=0.75cm] at (12.0, -1.6) {E/S};

    % rank 3: 10, 11, 12
    \node[compute, minimum width=3.6cm] at (0,    -2.4) {Resolvendo cenário 10};
    \node[io,      minimum width=0.75cm] at (3.6,  -2.4) {E/S};
    \node[compute, minimum width=3.6cm] at (4.35, -2.4) {Resolvendo cenário 11};
    \node[io,      minimum width=0.75cm] at (7.95, -2.4) {E/S};
    \node[compute, minimum width=3.7cm] at (8.7,  -2.4) {Resolvendo cenário 12};
    \node[io,      minimum width=0.75cm] at (12.4, -2.4) {E/S};

    % rank N: k, k+1, k+2
    \node[compute, minimum width=3.6cm] at (0,    -3.6) {Resolvendo cenário $k$};
    \node[io,      minimum width=0.75cm] at (3.6,  -3.6) {E/S};
    \node[compute, minimum width=3.5cm] at (4.35, -3.6) {Resolvendo cenário $k{+}1$};
    \node[io,      minimum width=0.75cm] at (7.85, -3.6) {E/S};
    \node[compute, minimum width=3.7cm] at (8.6,  -3.6) {Resolvendo cenário $k{+}2$};
    \node[io,      minimum width=0.75cm] at (12.3, -3.6) {E/S};

    %---- eixo do tempo ----
    \draw[timeaxis] (-0.1, -4.3) -- (13.2, -4.3);
    \node[font=\small, anchor=west] at (13.25, -4.3) {Tempo};

    %---- legenda ----
    \node[draw, rounded corners=2pt, fill=compcolor, minimum width=0.9cm, minimum height=0.45cm] at (1.0, -5.0) {};
    \node[font=\scriptsize, anchor=west] at (1.55, -5.0) {Computação --- resolução do cenário};
    \node[draw, rounded corners=2pt, fill=iocolor, minimum width=0.9cm, minimum height=0.45cm] at (7.5, -5.0) {};
    \node[font=\scriptsize, anchor=west] at (8.05, -5.0) {Escrita independente em \textit{offset} distinto};
  \end{tikzpicture}
  \caption{Linha do tempo da fase de simulação final do SDDP na arquitetura
  proposta. Os cenários $\{1, 2, \ldots\}$ são particionados entre os
  $N$ \textit{ranks} MPI: cada \textit{rank} resolve um subconjunto disjunto
  e, ao final da resolução de cada cenário, persiste seus resultados por meio
  de uma operação de E/S em um \textit{offset} distinto do(s) arquivo(s)
  compartilhado(s) via \texttt{MPI\_File\_write\_at}. Como a fase de simulação
  final é \textit{bag-of-tasks} (cenários mutuamente independentes) e as
  escritas não exigem coordenação coletiva, os \textit{ranks} executam suas
  operações de E/S de forma totalmente independente, sem sincronização entre
  processos.}
  \label{fig:timeline_sddp}
\end{figure}
